\chapter{The Common Framework}

Every program in this book begins by setting a few \texttt{Z} variables and
calling line 60000:

\begin{lstlisting}
ZT$="ROADHOG":ZA$="NOVA PORT STUB":ZS$="CROOKED ROAD DRIVING":ZP=1
GOSUB 60000
\end{lstlisting}

The common framework draws the opening screen, plays a short sound, waits for
the player, clears the machine state, and returns to the game.

\section{Framework Controls}

\begin{center}
\begin{tabularx}{0.92\textwidth}{>{\ttfamily}lX}
\toprule
ZT\$ & Title printed in the center of the splash screen.\\
ZA\$ & Author or program-series line.\\
ZS\$ & Short subtitle explaining the program.\\
ZP & Prompt mode. Use 0 for Enter or Space, 1 for Space start.\\
\bottomrule
\end{tabularx}
\end{center}

\section{What It Does}

The framework deliberately keeps to visible, beginner-friendly work:

\begin{itemize}
\item It clears Copper state and hides all hardware sprites.
\item It enters graphics mode and draws a high-color title frame.
\item It centers colorful bitmap titles with \texttt{GTEXT}, \texttt{LEN}, and \texttt{INT}.
\item It plays a short sound cue.
\item It waits for Enter or Space depending on \texttt{ZP}.
\item It returns to text mode and hands control back to the game.
\end{itemize}

\begin{novabox}
The framework is ordinary NovaBASIC. It is not a hidden ROM routine. Readers can
list it, change it, or replace it with their own house style.
\end{novabox}

\section{Why This Belongs in the Book}

A shared framework gives every program a little polish without making every
listing repeat the same title-screen work. It also lets the book teach a real
pattern: put common code in one place, keep its \texttt{Z} variables for
itself, and call it from the individual programs.

\begin{typeinbox}
Keep program-specific code out of \commonlisting. Sprites, maps, levels, and
special routines belong in the program that uses them, so each listing can be
studied and changed on its own.
\end{typeinbox}
